\section*{第 \ref{ch:b2:multint} 章——曲线积分与重积分}

\begin{solution}{exo:b2:multint:1}
(a) 线段 $\gamma(t) = (t, t)$，$t \in [0,1]$：
\[
\int_\gamma y^2\dd x + x\dd y
= \int_0^1 (t^2 + t)\,\dd t = \frac13 + \frac12 = \frac56 .
\]
(b) 抛物线 $\gamma(t) = (t, t^2)$：
\[
\int_0^1 \bigl(t^4\cdot 1 + t\cdot 2t\bigr)\dd t
= \frac15 + \frac23 = \frac{13}{15} .
\]
两个值不同，故该积分与路径有关：该形式
\emph{不}是\omterm{def:b2:multint:exact}{恰当形式}——与此相符，$P_y = 2y \neq 1 = Q_x$，
故它甚至不是\omterm{def:b2:multint:exact}{闭形式}。
\end{solution}

\begin{solution}{exo:b2:multint:2}
$P = 2xy + y^3$, $Q = x^2 + 3xy^2 + 1$: $P_y = 2x + 3y^2 = Q_x$,
在 $\R^2$ 上闭。求 $f$ 使 $f_x = P$：$f = x^2y + xy^3
+ g(y)$；再由 $f_y = x^2 + 3xy^2 + g'(y) = Q$ 得 $g'(y) = 1$，
取 $g(y) = y$。于是
\[
f(x, y) = x^2y + xy^3 + y
\]
即为一个\omterm{def:b2:multint:exact}{势函数}（$\R^2$ 是星形集，
故由庞加莱引理，\omterm{def:b2:multint:exact}{势函数}必然存在；直接写出更为简便）。由
\cref{thm:b2:multint:ftc}，对从 $(0,0)$ 到 $(1,2)$ 的任一曲线：
\[
\int_\gamma\omega = f(1, 2) - f(0, 0) = 2 + 8 + 2 = 12 .
\]
\end{solution}

\begin{solution}{exo:b2:multint:3}
积分区域为 $\{0 \leq x \leq 1,\ x^2 \leq y \leq x\}$。
先对 $y$ 积分：
\[
\int_0^1\!\int_{x^2}^{x}(x + y)\,\dd y\,\dd x
= \int_0^1\Bigl(x(x - x^2)
+ \frac{x^2 - x^4}{2}\Bigr)\dd x
= \int_0^1\Bigl(\frac{3x^2}{2} - x^3 - \frac{x^4}{2}\Bigr)\dd x
= \frac12 - \frac14 - \frac1{10} = \frac{3}{20} .
\]
先对 $x$ 积分：在高度 $y \in [0, 1]$ 处，截线为 $y \leq x \leq
\sqrt y$，于是
\[
\int_0^1\!\int_{y}^{\sqrt y}(x + y)\,\dd x\,\dd y
= \int_0^1\Bigl(\frac{y - y^2}{2}
+ y(\sqrt y - y)\Bigr)\dd y
= \frac14 - \frac16 + \frac25 - \frac13 = \frac{3}{20} .
\]
\end{solution}

\begin{solution}{exo:b2:multint:4}
\emph{第一个积分。} 在圆盘 $D_R$ 上取极坐标：
\[
\iint_{D_R}\frac{\dd x\,\dd y}{(1 + x^2 + y^2)^2}
= \int_0^{2\pi}\!\!\int_0^R
\frac{\rho\,\dd\rho\,\dd\alpha}{(1 + \rho^2)^2}
= 2\pi\Bigl[-\frac{1}{2(1 + \rho^2)}\Bigr]_0^R
= \pi\Bigl(1 - \frac{1}{1 + R^2}\Bigr)
\xrightarrow[R \to \infty]{} \pi .
\]
\emph{第二个积分。} 四分之一圆盘为 $0 \leq \alpha \leq
\frac\pi2$，$0 \leq \rho \leq 1$，且 $xy =
\rho^2\cos\alpha\sin\alpha$:
\[
\iint_{D'}xy\,\dd x\,\dd y
= \int_0^{\pi/2}\!\!\cos\alpha\sin\alpha\,\dd\alpha
\int_0^1 \rho^3\,\dd\rho
= \frac12\cdot\frac14 = \frac18 .
\]
\end{solution}

\begin{solution}{exo:b2:multint:5}
由\cref{cor:b2:multint:area}，取 $x = \cos^3 t$，$y = \sin^3 t$：
$x' = -3\cos^2 t\sin t$，$y' = 3\sin^2 t\cos t$，于是
\[
xy' - yx' = 3\cos^4 t\sin^2 t + 3\sin^4 t\cos^2 t
= 3\sin^2 t\cos^2 t = \frac{3}{4}\sin^2 2t
= \frac{3}{8}(1 - \cos 4t) .
\]
因此
\[
A = \frac12\int_0^{2\pi}\frac38(1 - \cos 4t)\,\dd t
= \frac{3}{16}\cdot 2\pi = \frac{3\pi}{8} .
\]
（\omterm{pb:b2:curves:1}{星形线}位于\omterm{def:b2:multint:domain}{面积}为 $\pi$ 的单位圆盘内；
$\pi$ 的八分之三对其四尖星形而言是合理的。）
\end{solution}

\begin{solution}{exo:b2:multint:6}
\emph{逐层累加：} 对每个 $(x, y)$（在单位圆盘 $D$ 内），$z$
从 $x^2 + y^2$ 变到 $1$：
\[
V = \iint_D \bigl(1 - x^2 - y^2\bigr)\dd x\,\dd y
= \int_0^{2\pi}\!\!\int_0^1 (1 - \rho^2)\rho\,\dd\rho\,\dd\alpha
= 2\pi\Bigl(\frac12 - \frac14\Bigr) = \frac\pi2 .
\]
\emph{截面法：} 高度 $z \in [0, 1]$ 处的截面为圆盘
$x^2 + y^2 \leq z$，其\omterm{def:b2:multint:domain}{面积}为 $\pi z$：
\[
V = \int_0^1 \pi z\,\dd z = \frac\pi2 .
\]
\end{solution}

\begin{solution}{exo:b2:multint:7}
在球坐标下，体积元为
$r^2\cos\varphi\,\dd r\,\dd\theta\,\dd\varphi$
（见\cref{ex:b2:multint:spherical}），角向部分
的积分为 $4\pi$（$2\pi$ 来自 $\theta$，$\int_{-\pi/2}^{\pi/2}\cos = 2$）。
球壳的体积为
\[
\int_a^b 4\pi r^2\,\dd r = \frac43\pi\bigl(b^3 - a^3\bigr).
\]
对第二个积分，被积函数 $1/r$ 只依赖于 $r$：
\[
\iiint_B \frac{\dd x\,\dd y\,\dd z}{r}
= \int_0^R 4\pi r^2\cdot\frac1r\,\dd r
= 4\pi\,\frac{R^2}{2} = 2\pi R^2 .
\]
（被积函数在原点处趋于无穷，但无碍：$r^2/r = r$
\omterm{def:b2:metric:continuity}{连续}——球壳 $\varepsilon \leq r
\leq R$ 上的积分当 $\varepsilon \to 0$ 时收敛，这正是该断言的精确含义。
此类计算是迈向牛顿定理的第一步，
该定理断言均匀球体的吸引作用
与位于其中心的质点相同。）
\end{solution}

\begin{solution}{exo:b2:multint:8}
被积函数 $g(t; x, y) = xP(tx, ty) + yQ(tx, ty)$ 为
$\mathcal{C}^1$ 类（关于 $(x, y)$），关于 $t$ \omterm{def:b2:metric:continuity}{连续}，偏
导数在 $[0,1] \times U$ 上\omterm{def:b2:metric:continuity}{连续}；由积分号下求导（见
\cref{ch:b2:integration}，应用于
$t$ 的紧区间 $[0,1]$ 上，此时控制条件自动成立）
可得
\[
f_x(x, y)
= \int_0^1 \bigl(P(tx, ty) + tx\,P_x(tx, ty)
+ ty\,Q_x(tx, ty)\bigr)\dd t .
\]
利用闭性 $Q_x = P_y$：
\[
tx\,P_x(tx, ty) + ty\,P_y(tx, ty)
= t\,\frac{\dd}{\dd t}\bigl[P(tx, ty)\bigr] ,
\]
于是被积函数为 $P(tx, ty) + t\frac{\dd}{\dd t}P(tx, ty) =
\frac{\dd}{\dd t}\bigl[t\,P(tx, ty)\bigr]$，从而
\[
f_x(x, y) = \Bigl[t\,P(tx, ty)\Bigr]_0^1 = P(x, y) .
\]
同理可得 $f_y = Q$（同样的计算，将 $P_y = Q_x$ 反向使用
）。注意条件在何处起作用：$f$ 由
沿线段 $[0, M]$ 积分定义，该线段位于 $U$ 中，
这正是因为 $U$ 为星形区域。
\end{solution}

\begin{solution}{exo:b2:multint:9}
在带形区域 $[0, A] \times [0, \infty)$ 上，函数 $(x, y)
\mapsto e^{-xy}\sin x$ 在 $y =
\infty$ 处并非朴素意义下一致绝对可积，但每个累次积分
都收敛，且二者相等可由 $[0, A]
\times [0, B]$ 上的富比尼定理及极限 $B \to \infty$ 得到（尾项
$\int_0^A\int_B^\infty e^{-xy}\abs{\sin x}\,\dd y\,\dd x \leq
\int_0^A \frac{e^{-Bx}\abs{\sin x}}{x}\dd x \leq \int_0^A e^{-Bx}
\dd x\to 0$，用到 $\abs{\sin x} \leq x$）。

\emph{先对 $y$ 积分：}$\int_0^\infty e^{-xy}\,\dd y = \frac1x$，当 $x
> 0$，故第一个积分为 $\int_0^A \frac{\sin x}{x}\,\dd x$。

\emph{先对 $x$ 积分：}两次分部积分（或取
$\int_0^A e^{(i - y)x}\dd x$ 的虚部）可得
\[
\int_0^A e^{-xy}\sin x\,\dd x
= \frac{1 - e^{-Ay}(y\sin A + \cos A)}{1 + y^2} .
\]
对 $y$ 在 $[0, \infty)$ 上积分，其中 $\int_0^\infty
\frac{\dd y}{1 + y^2} = \frac\pi2$ 单独分离出来：
\[
\int_0^A \frac{\sin x}{x}\,\dd x
= \frac{\pi}{2}
- \int_0^\infty e^{-Ay}\,\frac{y\sin A + \cos A}{1 + y^2}\,\dd y .
\]
余项的绝对值不超过 $\int_0^\infty e^{-Ay}\frac{y +
1}{1 + y^2}\dd y \leq \int_0^\infty e^{-Ay}\cdot\frac{1+y}{1+y^2}
\,\dd y \to 0$，当 $A \to \infty$ 时（控制收敛定理，或
粗略取界 $\frac{1 + y}{1 + y^2} \leq \frac32$，得
$\frac{3}{2A}$）。因此 $\int_0^\infty\frac{\sin x}{x}\dd x =
\frac\pi2$——与
在\cref{ch:b2:integration}中通过对参变量积分求导所得的值相同；
这里改由富比尼定理来完成。
\end{solution}

\begin{solution}{exo:b2:multint:10}
以弧长为参数 $s \in [0, 2\pi]$，于是 $x'^2 + y'^2 = 1$，
并平移坐标，使 $\int_0^{2\pi} x(s)\,\dd s = 0$。由
\cref{cor:b2:multint:area}，
\[
2A = \oint x\,\dd y - y\,\dd x
= \int_0^{2\pi}\bigl(xy' - yx'\bigr)\dd s .
\]
对 $\oint y\,\dd x$ 在一个周期上作分部积分（边界
项因周期性相消），$-\int yx' = \int y'x$，于是实际有
$2A = 2\int_0^{2\pi}xy'\,\dd s$。再由 $2xy' \leq x^2 + y'^2$
可得
\[
2A \leq \int_0^{2\pi}\bigl(x^2 + y'^2\bigr)\dd s
= \int_0^{2\pi} x^2 + \int_0^{2\pi}\bigl(1 - x'^2\bigr)
= 2\pi - \int_0^{2\pi}\bigl(x'^2 - x^2\bigr)\dd s .
\]
瓦廷格不等式（见\cref{ch:b2:fourier}中的练习：对
$2\pi$ 为周期、均值为零的 $\mathcal{C}^1$ 函数，
$\int x^2 \leq \int x'^2$）使最后一项的积分非负：
$A \leq \pi$。等号成立要求瓦廷格不等式中取等号（$x(s) =
a\cos s + b\sin s$），且 $2xy' \leq x^2 + y'^2$ 中取等号（$y' = x$
逐点成立），由此推出 $y = a\sin s - b\cos s + c$：该曲线为
单位圆（适当选取圆心）。由于长为 $L$ 的曲线
可伸缩为长 $2\pi$ 的曲线，因此一般结论为 $A \leq
\frac{L^2}{4\pi}$：在周长给定的所有闭曲线中，
圆所围\omterm{def:b2:multint:domain}{面积}最大。
\end{solution}

\begin{solution}{exo:b2:multint:11}
球坐标下 $z = r\sin\varphi$，且 $\dd x\,\dd
y\,\dd z = r^2\cos\varphi\,\dd r\,\dd\theta\,\dd\varphi$:
\[
\iiint_B z^2
= \int_0^R r^4\,\dd r\int_0^{2\pi}\dd\theta
\int_{-\pi/2}^{\pi/2}\sin^2\varphi\cos\varphi\,\dd\varphi
= \frac{R^5}5\cdot2\pi\cdot
\Bigl[\frac{\sin^3\varphi}3\Bigr]_{-\pi/2}^{\pi/2}
= \frac{4\pi R^5}{15}.
\]
由球体关于坐标置换的对称性，
$\iiint_B x^2 = \iiint_B y^2 = \iiint_B z^2$，故 $\iiint_B(x^2
+ y^2 + z^2) = 3\cdot\frac{4\pi R^5}{15} = \frac{4\pi R^5}5$.
壳层验算：$\int_0^R r^2\cdot 4\pi r^2\,\dd r = \frac{4\pi
R^5}5$——被积函数 $r^2$ 在半径为
$r$ 的球面上为常数，该球面\omterm{def:b2:multint:domain}{面积}为 $4\pi r^2$。
\end{solution}

\begin{solution}{exo:b2:multint:12}
取由 $(x_i, y_i)$ 到 $(x_{i+1}, y_{i+1})$ 的边，
参数化为 $\gamma(t) = \bigl((1-t)x_i + tx_{i+1},\ (1-t)y_i +
ty_{i+1}\bigr)$。该边对
$\frac12\oint(x\,\dd y - y\,\dd x)$ 的贡献为
\[
\frac12\int_0^1\Bigl(\bigl((1-t)x_i +
tx_{i+1}\bigr)(y_{i+1} - y_i) - \bigl((1-t)y_i +
ty_{i+1}\bigr)(x_{i+1} - x_i)\Bigr)\dd t ,
\]
由于 $\int_0^1\bigl((1-t)u + tv\bigr)\dd t = \frac{u +
v}2$，此式等于
\[
\frac14\Bigl((x_i + x_{i+1})(y_{i+1} - y_i) - (y_i +
y_{i+1})(x_{i+1} - x_i)\Bigr)
= \frac12\bigl(x_iy_{i+1} - x_{i+1}y_i\bigr),
\]
交叉项相消。对 $m$ 条边求和，
即得鞋带公式（见\cref{cor:b2:multint:area}）。
三角形 $(0,0), (1,0), (0,1)$：$\frac12\bigl((0\cdot0 - 1\cdot0) +
(1\cdot1 - 0\cdot0) + (0\cdot0 - 0\cdot1)\bigr) = \frac12$，
即正确的\omterm{def:b2:multint:domain}{面积}。
\end{solution}

\begin{solution}{pb:b2:multint:1}
\textbf{1.} 球 $B_n(R)$ 由逐次界描述
$-R \leq x_n \leq R$，然后 $\abs{x_{n-1}} \leq \sqrt{R^2 -
x_n^2}$，依此类推；将 $x_i = Ru_i$ 代入
$n$ 个一元积分中，每个都乘以 $R$，并把
积分限映为 $B_n(1)$ 的积分限：$v_n(R) = R^n\,v_n(1) =
V_nR^n$。

\textbf{2.} 沿 $x_n = t$ 截取：$B_n(1)$ 的截面为
球 $B_{n-1}\bigl(\sqrt{1 - t^2}\bigr)$，故由
第 1 问，
\[
V_n = \int_{-1}^1 v_{n-1}\bigl(\sqrt{1 - t^2}\bigr)\,\dd t
= V_{n-1}\int_{-1}^{1}(1 - t^2)^{\frac{n-1}2}\,\dd t .
\]

\textbf{3.} 令 $t = \sin\theta$，$\dd t =
\cos\theta\,\dd\theta$，且 $(1 - t^2)^{\frac{n-1}2} =
\cos^{n-1}\theta$，均在 $\intcc{-\pi/2}{\pi/2}$ 上：
\[
\int_{-1}^1(1 - t^2)^{\frac{n-1}2}\dd t
= \int_{-\pi/2}^{\pi/2}\cos^n\theta\,\dd\theta
= 2\int_0^{\pi/2}\cos^n\theta\,\dd\theta = 2W_n,
\]
由 $\theta \mapsto \frac\pi2 - \theta$ 可
互换 $W_n$ 的正弦形式与余弦形式。

\textbf{4.} 写 $\sin^n = \sin^{n-2}(1 - \cos^2)$，并
对 $\int\sin^{n-2}\cos\cdot\cos$ 分部积分（取 $v =
\frac{\sin^{n-1}}{n-1}$）：
\[
W_n = W_{n-2} - \frac{W_n}{n-1}
\quad\Longrightarrow\quad
W_n = \frac{n-1}{n}W_{n-2}.
\]
因此 $nW_nW_{n-1} = (n-1)W_{n-1}W_{n-2}$：数列
$(nW_nW_{n-1})$ 为常数，等于 $1\cdot W_1W_0 =
1\cdot\frac\pi2$，故 $W_nW_{n-1} = \frac{\pi}{2n}$。

\textbf{5.} 由第 2--3 问得 $V_n = 2W_nV_{n-1}$，连用两次：
\[
V_n = 2W_n\cdot 2W_{n-1}\,V_{n-2}
= 4\,\frac{\pi}{2n}\,V_{n-2} = \frac{2\pi}n\,V_{n-2}.
\]

\textbf{6.} 由 $V_0 = 1$，得 $V_{2k} = \frac{2\pi}{2k}V_{2k-2}
= \frac\pi kV_{2k-2}$，故 $V_{2k} = \frac{\pi^k}{k!}$，由
归纳法可得。由 $V_1 = 2$：$V_{2k+1} =
\frac{2\pi}{2k+1}V_{2k-1}$，故
\[
V_{2k+1} = 2\prod_{j=1}^k\frac{2\pi}{2j+1}
= \frac{2^{k+1}\pi^k}{1\cdot3\cdot5\cdots(2k+1)} .
\]

\textbf{7.} $V_1 = 2$，$V_2 = \pi \approx 3.142$，$V_3 =
\frac{4\pi}3 \approx 4.189$，$V_4 = \frac{\pi^2}2 \approx
4.935$，$V_5 = \frac{8\pi^2}{15} \approx 5.264$，$V_6 =
\frac{\pi^3}6 \approx 5.168$，$V_7 = \frac{16\pi^3}{105}
\approx 4.725$。相邻项之比为 $V_n/V_{n-1} = 2W_n$，
且 $(W_n)$ 单调递减（$\sin^n \leq \sin^{n-1}$
逐点成立）。现 $2W_5 = 2\cdot\frac45\cdot\frac23 =
\frac{16}{15} > 1$，而 $2W_6 =
2\cdot\frac56\cdot\frac34\cdot\frac12\cdot\frac\pi2 =
\frac{5\pi}{16} < 1$：比值大于 $1$ 直至 $n = 5$，而
比值小于 $1$ 始于 $n = 6$——$(V_n)$ 先增，
至其最大值 $V_5$，随后递减。

\textbf{8.} 当 $n \geq 13 > 4\pi$ 时：$V_n/V_{n-2} = 2\pi/n <
\tfrac12$，故 $V_{n} \leq C\cdot 2^{-n/2}$，其中
常数固定；更好地，对任意 $q > 0$，$2\pi/n < q^2$ 当
$n$ 充分大时，故 $V_n/q^n \to 0$：此衰减快于任何
等比数列。级数 $\sum V_n$ 的收敛性由比值
$V_n/V_{n-2} \to 0$ 推出（自
某项起与等比级数比较）。

\textbf{9.} $\sum_{k\geq0}V_{2k}x^{2k} =
\sum_{k\geq0}\frac{(\pi x^2)^k}{k!} = \eu^{\pi x^2}$，即
指数级数（见\cref{ch:b2:powerseries}），对
每个 $x$ 收敛。取 $x = 1$：$\sum_kV_{2k} = \eu^\pi \approx
23.14$。

\textbf{10.} 在立方体 $\intcc{-R}R^n$ 上，被积函数为
乘积 $\prod_i\eu^{-x_i^2}$，故累次积分
可分解为 $\bigl(\int_{-R}^R\eu^{-t^2}\dd t\bigr)^n$。
令 $R \to \infty$，并利用
$\int_\R\eu^{-t^2}\dd t = \sqrt\pi$
（见\cref{ex:b2:multint:polar}）：$I_n = \pi^{n/2}$。

\textbf{11.} 令 $t = u^2$，得 $\Gamma(\tfrac12) =
\int_0^\infty t^{-1/2}\eu^{-t}\dd t =
2\int_0^\infty\eu^{-u^2}\dd u = \sqrt\pi$。反复利用
$\Gamma(s+1) = s\Gamma(s)$: $\Gamma(k+1) = k!\,\Gamma(1) =
k!$，并有
\[
\Gamma\Bigl(k + \frac32\Bigr)
= \Bigl(k + \frac12\Bigr)\Bigl(k - \frac12\Bigr)\cdots
\frac12\cdot\Gamma\Bigl(\frac12\Bigr)
= \frac{(2k+1)(2k-1)\cdots1}{2^{k+1}}\,\sqrt\pi .
\]

\textbf{12.} 设 $F_n = \pi^{n/2}/\Gamma(\frac n2 + 1)$。
由于 $\Gamma(\frac n2 + 1) = \frac n2\,\Gamma(\frac n2) =
\frac n2\,\Gamma(\frac{n-2}2 + 1)$，于是 $F_n =
\frac{2\pi}nF_{n-2}$，与 $V_n$ 的递推相同（第
5 题）。初值为 $F_1 = \sqrt\pi/\Gamma(\frac32) =
\sqrt\pi/(\frac{\sqrt\pi}2) = 2 = V_1$ 及 $F_2 =
\pi/\Gamma(2) = \pi = V_2$。由归纳法可得 $V_n = F_n$，对
一切 $n$ 成立；第 11 题把它还原为
第 6 题的两种闭式。

\textbf{13.} 令 $r = \sqrt t$，则 $\int_0^\infty
\eu^{-r^2}r^{n-1}\dd r = \frac12\int_0^\infty
t^{\frac n2 - 1}\eu^{-t}\dd t = \frac12\Gamma(\frac n2)$.
于是
\[
n\,V_n\int_0^\infty\eu^{-r^2}r^{n-1}\dd r
= V_n\cdot\frac n2\,\Gamma\Bigl(\frac n2\Bigr)
= V_n\,\Gamma\Bigl(\frac n2 + 1\Bigr) = \pi^{n/2} = I_n .
\]
两边均已证明，该恒等式可以\emph{理解}为
高斯积分的壳层分解：半径为 $r$ 的球面
\omterm{def:b2:multint:domain}{面积}为 $nV_nr^{n-1}$，而高斯
权重 $\eu^{-r^2}$ 沿各壳层积分。

\textbf{14.} $s_0 = V_1 = 2$（$0$ 维球面是两点），
$s_1 = 2V_2 = 2\pi$, $s_2 = 3V_3 = 4\pi$, $s_3 = 4V_4 =
2\pi^2$，并有 $\int_0^R s_{n-1}r^{n-1}\dd r = V_nR^n =
v_n(R)$：\omterm{def:b2:multint:domain}{面积}即体积关于半径的导数。

\textbf{15.} 设 $n = 2k$，由斯特林公式
（见\cref{thm:b2:comparison:stirling}）得 $k! \sim
\sqrt{2\pi k}\,(k/\eu)^k$，于是
\[
V_{2k} = \frac{\pi^k}{k!}
\sim \frac{(\pi\eu/k)^k}{\sqrt{2\pi k}}
= \frac1{\sqrt{\pi n}}\Bigl(\frac{2\pi\eu}n\Bigr)^{n/2}
\qquad (n = 2k).
\]
对奇数 $n$：$V_{2k+1} = 2W_{2k+1}V_{2k} \leq 2V_{2k}$，于是
比几何衰减更快的界同样成立（相差一个
因子 $2$ 且指数平移一位）——对每个 $q >
0$，$V_n = o(q^n)$。

\textbf{16.} $V_2/4 = \pi/4 \approx 0.785$；$V_3/8 = \pi/6
\approx 0.524$; $V_{10}/2^{10} = \frac{\pi^5}{120\cdot1024}
\approx 0.0025$。一般地 $\frac{V_n/2^n}{V_{n-2}/2^{n-2}}
= \frac{2\pi}{4n} = \frac{\pi}{2n} \to 0$：比值趋于
$0$（比几何衰减更快）。内切球所占比例
趋于零：立方体的体积集中到
其顶点附近。

\textbf{17.} 由第 1 题，半径为 $1 -
\varepsilon$ 的内球体积为 $V_n(1-\varepsilon)^n$，于是外层
壳层所占比例为 $1 - (1 - \varepsilon)^n \to 1$。
当 $n = 100$，$\varepsilon = 0.01$ 时：$(0.99)^{100} =
\eu^{100\ln0.99} \approx \eu^{-1.005} \approx 0.366$：约
$63\%$ 的球体位于距其表面 $1\%$ 以内。

\textbf{18.} $(W_n)$ 递减，故 $W_n \leq W_{n-1} \leq
W_{n-2} = \frac{n}{n-1}W_n$：夹逼，$W_{n-1}/W_n \to 1$。
乘以 $W_nW_{n-1} = \frac\pi{2n}$：$W_n^2 \sim
\frac\pi{2n}$，即 $W_n \sim \sqrt{\pi/(2n)}$。下界：
$W_n^2 \geq W_nW_{n+1} = \frac{\pi}{2(n+1)}$，故 $W_n \geq
\sqrt{\pi/(2(n+1))}$ 对每个 $n$ 成立。

\textbf{19.} 分子：由 $1 - x^2 \leq \eu^{-x^2}$ 得 $(1 -
x^2)^{\frac{n-1}2} \leq \eu^{-(n-1)x^2/2}$，并令 $a =
\frac{n-1}2$,
\[
\int_\delta^1(1 - x^2)^{\frac{n-1}2}\dd x
\leq \int_\delta^\infty\eu^{-ax^2}\dd x
\leq \int_\delta^\infty\frac x\delta\,\eu^{-ax^2}\dd x
= \frac{\eu^{-a\delta^2}}{2a\delta}
= \frac{\eu^{-(n-1)\delta^2/2}}{(n-1)\delta}.
\]
分母：由第 18 题，$2W_n \geq \sqrt{2\pi/(n+1)}$。
相除即得上面的界。取 $\delta =
s/\sqrt{n-1}$，此时上式化为
$\sqrt{\tfrac{n+1}{2\pi(n-1)}}\;\eu^{-s^2/2}/s = O\bigl(
\eu^{-s^2/2}/s\bigr)$，关于 $n$ 一致：在带状区域
$\abs{x_1} \leq s/\sqrt{n-1}$ 之外几乎没有体积，当
$s$ 适当大时——由对称性，对
每个方向同样成立。

\textbf{20.} 两个结论并不矛盾，因为它们刻画的是
同一点的不同侧面。几乎每个点
$B_n(1)$ 的\omterm{def:b2:nvs:norm}{范数}接近 $1$（第 17 题：径向
集中于球面附近），然而它的 $n$ 个
坐标都很小，量级为 $1/\sqrt n$（第 19 题），
这与 $n$ 个大小为 $1/\sqrt n$ 的坐标
\omterm{def:b2:nvs:norm}{范数}量级为 $1$ 是相容的。高维体积集中在
所有坐标均分\omterm{def:b2:nvs:norm}{范数}预算之处——
靠近球面，但远离每个坐标轴的极点。

\textbf{21.} 取 $\Delta_n$ 在 $x_n = t \in \intcc01$ 处的
截面为 $\{x' \in \R^{n-1} : x_i \geq 0,\ \sum x_i \leq 1 -
t\} = (1-t)\Delta_{n-1}$，体积为
$(1-t)^{n-1}\operatorname{vol}(\Delta_{n-1})$，由齐次性。
于是
\[
\operatorname{vol}(\Delta_n) =
\operatorname{vol}(\Delta_{n-1})\int_0^1(1 - t)^{n-1}\dd t =
\frac{\operatorname{vol}(\Delta_{n-1})}{n}
\quad\Longrightarrow\quad
\operatorname{vol}(\Delta_n) = \frac1{n!}\,.
\]

\textbf{22.} $2^n$ 个符号卦限把 $C_n$ 切分为 $2^n$
个 $\Delta_n$ 的副本（在坐标
超平面的重叠可忽略）：$\operatorname{vol}(C_n) = \frac{2^n}{n!}$。
若 $\sum\abs{x_i} \leq 1$，则 $\sum x_i^2 \leq
\bigl(\sum\abs{x_i}\bigr)^2 \leq 1$，故 $C_n \subseteq B_n(1)$；
又 $B_n(1) \subseteq \intcc{-1}1^n$，因为 $\abs{x_i} \leq
\norm x$。于是 $\frac{2^n}{n!} \leq V_n \leq 2^n$——
与第 15 题一致，它把 $V_n$ 置于
阶乘尺度与几何尺度之间。

\textbf{23.} 对单位圆盘内的 $(x, y)$，
$B_4(1)$ 在该处的截面为半径 $\sqrt{1 - x^2 - y^2}$ 的圆盘，
位于 $(z, w)$ 平面上，\omterm{def:b2:multint:domain}{面积}为 $\pi(1 - x^2 - y^2)$。用极坐标
表示，有：
\[
V_4 = \iint_{x^2+y^2\leq1}\pi(1 - x^2 - y^2)\,\dd x\,\dd y
= \pi\int_0^{2\pi}\!\!\int_0^1(1 - \rho^2)\rho\,
\dd\rho\,\dd\alpha
= \pi\cdot2\pi\cdot\frac14 = \frac{\pi^2}2 ,
\]
与第 6 题结果相符。

\textbf{24.} 所求概率即体积之比
$\dfrac{V_{20}}{2^{20}} = \dfrac{\pi^{10}}{10!\cdot2^{20}}
\approx \dfrac{0.0258}{1\,048\,576} \approx
2.5\cdot10^{-8}$。首次命中前的抽样次数
约为其倒数的量级，即 $4\cdot10^7$：拒绝抽样法
在圆盘上效果极佳（命中率为 $\pi/4$），此法
在 $20$ 维中毫无用处——维数灾难
在一行之内便清晰可见。

\textbf{25.} 第一条路线（第一至二部分）用到：
对累次积分作富比尼型切片，在每个坐标上作单变量
代换（利用齐次性），以及\omterm{pb:b2:multint:1}{沃利斯积分}
——纯粹的一元微积分加归纳法。第二条路线
（第三部分）用到：$I_n$ 乘积结构上的富比尼定理，
借助高斯积分作极坐标换元（
见\cref{ex:b2:multint:polar}），以及 $\Gamma$ 函数的
函数方程。两者汇合于 $V_n =
\pi^{n/2}/\Gamma(\frac n2 + 1)$，再借助斯特林公式
（见\cref{thm:b2:comparison:stirling}）将公式化为渐近式。
本科三年级分册将以上全部内容重新建立在
勒贝格积分之上。在那里，富比尼定理与变量代换
对一般可积函数都是定理，球
坐标在每个维数中都存在，同样的球体积
重新出现，作为乘积测度与斯特林问题的
现成的算例——用控制
收敛定理取代我们手工构造的夹逼。
\end{solution}
